\subsubsection{Movimiento en coordenadas cilíndricas}

Se observa una partícula en movimiento con respecto a un sistema de referencia inercial. La trayectoria está dada por las siguientes funciones:
\[ r = Ae^{k\theta}, \quad z = hr \]
donde $r, \theta$ y $z$ son coordenadas cilíndricas, siendo $A, k$ y $h$ valores escalares constantes y mayores que cero, como se muestra en la Figura \ref{fig:c1_e01_p02_espiral_cilindricas}. Suponiendo que su rapidez es constante ($v_0$) y conocida:

\begin{enumerate}
    \item Calcule la velocidad $\vec{v}$ de la partícula en función de $\theta, A, k, h$ y $v_0$.
    \item Encuentre su aceleración $\vec{a}$ en función de los mismos parámetros.
    \item Pruebe que $\vec{a} \perp \vec{v}$.
    \item Encuentre una expresión para $\theta(t)$.
\end{enumerate}

% Figura principal
\begin{figure}[!ht]
    \centering
    \begin{tikzpicture}[>=Latex, scale=1.5]

        % --- CONFIGURACIÓN DE PARÁMETROS ---
        \def\A{0.15}      % Constante A de la espiral [u]
        \def\k{0.22}      % Constante k (crecimiento) [-]
        \def\thP{560}     % Ángulo del punto P [deg] (una vuelta y media)
        \def\thPmod{200}  % Azimut equivalente de P [deg]

        % --- SISTEMAS DE COORDENADAS ---
        \draw[->, thick, gris1] (-2.2,0) -- (2.4,0) node[below right] {$x$};
        \draw[->, thick, gris1] (0,-1.8) -- (0,2.2) node[left] {$y$};

        % --- GEOMETRÍA PRINCIPAL ---
        % Espiral logarítmica r = A e^{k theta} (vista superior)
        \draw[thick, gris1, domain=0:720, samples=300, smooth, variable=\t]
            plot ({\t}:{\A*exp(\k*\t/57.29578)});

        % --- PUNTOS Y COORDENADAS ---
        \pgfmathsetmacro{\rP}{\A*exp(\k*\thP/57.29578)}
        \coordinate (O) at (0,0);
        \coordinate (P) at (\thP:\rP);
        \draw[mydashed, gris1!60] (O) -- (P);

        % --- VECTORES Y FUERZAS ---
        % Vectores unitarios cilíndricos en P
        \draw[->, very thick, myred] (P) -- ++(\thP:0.55) node[above left] {$\hat{e}_r$};
        \draw[->, very thick, myred] (P) -- ++(\thP+90:0.55) node[below left] {$\hat{e}_\theta$};
        % Velocidad (tangente a la espiral, rapidez constante v0)
        \pgfmathsetmacro{\alphat}{atan(1/\k)}
        \draw[->, very thick, azul] (P) -- ++(\thP+\alphat:0.9) node[below] {$\vec{v}$};

        % --- ANOTACIONES Y ETIQUETAS ---
        \fill[gris1] (P) circle (1.2pt) node[above right, xshift=1pt] {$P$};
        \node[below left] at (O) {$O$};
        \draw[->] (0.35,0) arc (0:\thPmod:0.35);
        \node at ({\thPmod/2}:0.55) {$\theta$};
        \node[gris1] at (1.6,-1.4) {$r = Ae^{k\theta}$};

    \end{tikzpicture}
    \caption{Vista superior de la espiral $r=Ae^{k\theta}$; la altura crece como $z=hr$}
    \label{fig:c1_e01_p02_espiral_cilindricas}
\end{figure}

